%% maestro2tex -- generated figure; do not edit by hand.
\providecommand{\MaestroRoot}{include/analog/}
\begin{figure}[htbp]
  \centering
  {\pgfplotsset{compat=1.18}%
  \begin{tikzpicture}
    \begin{axis}[
      width=0.82\linewidth, height=6.2cm,
      xlabel={Test-source current $i_{\mathrm{we}}$ drawn out of the WE node~[\si{\micro\ampere}]}, ylabel={$V_{\mathrm{WE}}-v_{\mathrm{cm}}$~[\si{\milli\volt}]},
      grid=both,
      major grid style={line width=0.2pt, draw=black!14},
      minor grid style={line width=0.2pt, draw=black!7},
      minor tick num=1,
      axis line style={line width=0.4pt, draw=black!45},
      tick style={line width=0.4pt, draw=black!45},
      tick align=outside,
      label style={font=\small}, tick label style={font=\small},
      enlarge x limits=0.02, enlarge y limits=0.10,
      every axis plot/.append style={line join=round, no marks},
      legend style={font=\footnotesize, draw=none, fill=none,
                    at={(0.5,1.02)}, anchor=south, legend columns=-1,
                    /tikz/every even column/.append style={column sep=9pt}},
      legend cell align=left,
      scaled y ticks=false, scaled x ticks=false,
      y tick label style={/pgf/number format/fixed,
                          /pgf/number format/precision=1}
    ]
      \addplot[mtxA, line width=1.1pt, solid] table {\MaestroRoot data/tia_weerr_00.dat};
      \addlegendentry{tt}
      \addplot[mtxB, line width=1.1pt, dashed] table {\MaestroRoot data/tia_weerr_01.dat};
      \addlegendentry{ff}
      \addplot[mtxC, line width=1.1pt, dotted] table {\MaestroRoot data/tia_weerr_02.dat};
      \addlegendentry{fs}
      \addplot[mtxD, line width=1.1pt, dashdotted] table {\MaestroRoot data/tia_weerr_03.dat};
      \addlegendentry{sf}
      \addplot[mtxE, line width=1.1pt, densely dashed] table {\MaestroRoot data/tia_weerr_04.dat};
      \addlegendentry{ss}
    \end{axis}
  \end{tikzpicture}}
  \caption[{Error between the working-electrode potential and $v_{\mathrm{cm}}$ across the same sweep --- the quality of\,\dots}]{Error between the working-electrode potential and $v_{\mathrm{cm}}$ across the same sweep --- the quality of the virtual ground the electrode sees. It stays inside \SI{\pm 0.4}{\milli\volt} over the full current range in both directions, and its slope is the input impedance of the stage, set by the loop gain available at the summing node. Against a \SI{1.25}{\volt} full-scale output this is under \SI{0.04}{\percent} of full scale, so the virtual ground is not an accuracy limit at this gain setting.}
  \label{fig:tia-weerr}
\end{figure}
